\textbf{1. Kiến thức cần nhớ}
\begin{itemize}
\item Phương pháp giải bài toán bằng cách lập phương trình bậc hai tương tự như lập hệ phương trình, nhưng chỉ chọn \textbf{một ẩn số} chính, các đại lượng khác đều được biểu diễn theo ẩn số đó dẫn đến phương trình bậc hai $ax^2 + bx + c = 0$.
\end{itemize}

\textbf{2. Phương pháp giải các dạng}

%% DANG 9D6?21-1
\textbf{Dạng 1. Toán chuyển động}
\begin{enumerate}
\item Gọi vận tốc lúc đi (hoặc vận tốc dự định) là $x$ ($x > 0$).
\item Biểu diễn thời gian đi, thời gian về theo quãng đường và vận tốc.
\item Lập phương trình từ chênh lệch thời gian: $t_1 - t_2 = \Delta t$.
\item Quy đồng khử mẫu đưa về phương trình bậc hai, giải tìm $x$ và đối chiếu điều kiện $x > 0$.
\end{enumerate}

%% DANG 9D6?21-2
\textbf{Dạng 2. Toán năng suất, làm chung – làm riêng}
\begin{enumerate}
\item Gọi thời gian làm riêng một mình hoàn thành công việc là $x$ (ngày/giờ).
\item Năng suất làm trong $1$ ngày là $\dfrac{1}{x}$.
\item Lập phương trình theo tổng năng suất hoặc phần việc làm được đưa về phương trình bậc hai.
\end{enumerate}

%% DANG 9D6?21-3
\textbf{Dạng 3. Toán hình học (diện tích, kích thước)}
\begin{enumerate}
\item Gọi chiều rộng (hoặc chiều dài) của hình chữ nhật là $x$.
\item Biểu diễn kích thước còn lại và lập phương trình diện tích: $x(k \pm x) = S$.
\item Bài toán mở rộng lối đi, làm đường bao quanh vườn: $(a + 2x)(b + 2x) = S_{\text{mới}}$.
\item Định lí Pythagore trong tam giác vuông: $x^2 + (x + d)^2 = c^2$.
\end{enumerate}

%% DANG 9D6?21-4
\textbf{Dạng 4. Toán thực tế khác (tăng trưởng, kinh tế)}
\begin{enumerate}
\item Lãi suất kép / Tăng trưởng dân số sau $2$ năm: $A = A_0(1 + r\%)^2$.
\item Doanh thu = Số lượng bán $\times$ Đơn giá (với số lượng giảm khi giá tăng).
\end{enumerate}
